% GENERATED FILE. Edit canonical lessons, then run npm run generate:books.
% canonical-source-hash: fnv1a64:d12b8d3ce827b9b3
% canonical-lessons: RU-C332-zapekat, RU-C332-protiven, RU-C332-razogrevat, RU-C332-podgorat, RU-C332-obzharivat

\chapter{In the oven}
\label{ch:ru-in-the-oven}
\hlchaptermodality{332}

\begin{chapteropening}
\textbf{By the end of this chapter:} \emph{I can describe oven cooking: baking, a baking tray, heating up, food catching and browning.}

Complete the last lesson of chapter 332: I can describe oven cooking: baking, a baking tray, heating up, food catching and browning.
\end{chapteropening}

\section[\texorpdfstring{\ru{запекать} (zapekát)}{запекать (zapekát)}]{\ru{запекать} (zapekát) — to bake, to roast}

\label{lesson:RU-C332-zapekat}



\begin{quote}
Before the new one: say the Russian for to shape, then the Russian for a filling.
\end{quote}

\subsection*{You'll want to know: \ru{запекать}}
\textbf{\ru{запекать}} (\emph{zapekát}) — \textquotedblleft{}to bake, to roast\textquotedblright{}.

\begin{quote}
\ru{Мама} \ru{запекает} \ru{рыбу} \ru{в} \ru{духовке}.
\end{quote}

\begin{quote}
\emph{Máma zapekáyet rýbu v dukhóvke.}
\end{quote}

\begin{quote}
Mum is baking fish in the oven.
\end{quote}

Imperfective; perfective \emph{\ru{запечь}} (\emph{zapéch}), \emph{\ru{запеку}}, \emph{\ru{запечёшь}} (\emph{zapechyósh}).

\begin{cousinweb}
\emph{\ru{запекать}} is built on \emph{\ru{за}-} and \emph{\ru{печь}} \textquotedblleft{}to bake\textquotedblright{}, an inherited Slavic verb (Proto-Slavic \emph{pekti}); EtymDB lists no parent for the compound.
\end{cousinweb}

\begin{grammarlens}[title={What you've built}]
One more word for telling, explaining and describing at B1.
\end{grammarlens}

\subsection*{Your turn}
\begin{itemize}
\raggedright
  \item \emph{Say it:} \emph{zapekát}
  \item \emph{Say it:} \emph{zapekát}, once more
  \item \emph{Say it:} a sentence of your own with \emph{zapekát}
  \item \emph{Recall:} say the Russian for to shape, then the Russian for a filling, then say \emph{zapekát} again
\end{itemize}

\begin{tcolorbox}[breakable,colback=teal!4,colframe=teal!35!black,title={Before you move on}]
What does \ru{запекать} mean? (\textquotedblleft{}To bake, to roast\textquotedblright{}.) Say it once more.
\end{tcolorbox}
\section[\texorpdfstring{\ru{противень} (prótiven)}{противень (prótiven)}]{\ru{противень} (prótiven) — a baking tray}

\label{lesson:RU-C332-protiven}



\begin{quote}
Before the new one: say the Russian for a filling, then the Russian for to bake.
\end{quote}

\subsection*{You'll want to know: \ru{противень}}
\textbf{\ru{противень}} (\emph{prótiven}) — \textquotedblleft{}a baking tray\textquotedblright{}.

\begin{quote}
\ru{Печенье} \ru{лежит} \ru{на} \ru{противне}.
\end{quote}

\begin{quote}
\emph{Pechénye lezhít na prótivne.}
\end{quote}

\begin{quote}
The biscuits are on the baking tray.
\end{quote}

Masculine; genitive \emph{\ru{противня}}.

\begin{cousinweb}
\emph{\ru{противень}} is usually traced to \emph{\ru{против}} \textquotedblleft{}opposite\textquotedblright{}, a tray set facing the fire; EtymDB lists no parent.
\end{cousinweb}

\begin{grammarlens}[title={What you've built}]
One more word for telling, explaining and describing at B1.
\end{grammarlens}

\subsection*{Your turn}
\begin{itemize}
\raggedright
  \item \emph{Say it:} \emph{prótiven}
  \item \emph{Say it:} \emph{prótiven}, once more
  \item \emph{Say it:} a sentence of your own with \emph{prótiven}
  \item \emph{Recall:} say the Russian for a filling, then the Russian for to bake, then say \emph{prótiven} again
\end{itemize}

\begin{tcolorbox}[breakable,colback=teal!4,colframe=teal!35!black,title={Before you move on}]
What does \ru{противень} mean? (\textquotedblleft{}A baking tray\textquotedblright{}.) Say it once more.
\end{tcolorbox}
\section[\texorpdfstring{\ru{разогревать} (razogrevát)}{разогревать (razogrevát)}]{\ru{разогревать} (razogrevát) — to heat up, to warm up}

\label{lesson:RU-C332-razogrevat}



\begin{quote}
Before the new one: say the Russian for to bake, then the Russian for a baking tray.
\end{quote}

\subsection*{You'll want to know: \ru{разогревать}}
\textbf{\ru{разогревать}} (\emph{razogrevát}) — \textquotedblleft{}to heat up, to warm up\textquotedblright{}.

\begin{quote}
\ru{Я} \ru{разогрел} \ru{суп} \ru{на} \ru{плите}.
\end{quote}

\begin{quote}
\emph{Ya razogrél sup na plité.}
\end{quote}

\begin{quote}
I heated up the soup on the stove.
\end{quote}

Imperfective; perfective \emph{\ru{разогреть}} (\emph{razogrét}), \emph{\ru{разогрею}}.

\begin{cousinweb}
\emph{\ru{разогревать}} is built on \emph{\ru{раз}-}, \emph{\ru{о}-} and \emph{\ru{греть}} \textquotedblleft{}to warm\textquotedblright{}, inherited from Proto-Slavic \emph{grěti} \textquotedblleft{}to heat, to warm\textquotedblright{}, from PIE \emph{g\textsuperscript{w}\textsuperscript{h}er-} (EtymDB).
\end{cousinweb}

\begin{grammarlens}[title={What you've built}]
One more word for telling, explaining and describing at B1.
\end{grammarlens}

\subsection*{Your turn}
\begin{itemize}
\raggedright
  \item \emph{Say it:} \emph{razogrevát}
  \item \emph{Say it:} \emph{razogrevát}, once more
  \item \emph{Say it:} a sentence of your own with \emph{razogrevát}
  \item \emph{Recall:} say the Russian for to bake, then the Russian for a baking tray, then say \emph{razogrevát} again
\end{itemize}

\begin{tcolorbox}[breakable,colback=teal!4,colframe=teal!35!black,title={Before you move on}]
What does \ru{разогревать} mean? (\textquotedblleft{}To heat up, to warm up\textquotedblright{}.) Say it once more.
\end{tcolorbox}
\section[\texorpdfstring{\ru{подгорать} (podgorát)}{подгорать (podgorát)}]{\ru{подгорать} (podgorát) — to get burnt (of food), to catch}

\label{lesson:RU-C332-podgorat}



\begin{quote}
Before the new one: say the Russian for a baking tray, then the Russian for to heat up.
\end{quote}

\subsection*{You'll want to know: \ru{подгорать}}
\textbf{\ru{подгорать}} (\emph{podgorát}) — \textquotedblleft{}to get burnt (of food), to catch\textquotedblright{}.

\begin{quote}
\ru{Пирог} \ru{немного} \ru{подгорел}.
\end{quote}

\begin{quote}
\emph{Piróg nemnógo podgorél.}
\end{quote}

\begin{quote}
The pie got a little burnt.
\end{quote}

Imperfective; perfective \emph{\ru{подгореть}} (\emph{podgorét}).

\begin{cousinweb}
\emph{\ru{подгорать}} is built on \emph{\ru{под}-} \textquotedblleft{}underneath\textquotedblright{} and \emph{\ru{гореть}} \textquotedblleft{}to burn\textquotedblright{}, inherited from Proto-Slavic \emph{gorěti} (EtymDB): to burn at the bottom.
\end{cousinweb}

\begin{grammarlens}[title={What you've built}]
One more word for telling, explaining and describing at B1.
\end{grammarlens}

\subsection*{Your turn}
\begin{itemize}
\raggedright
  \item \emph{Say it:} \emph{podgorát}
  \item \emph{Say it:} \emph{podgorát}, once more
  \item \emph{Say it:} a sentence of your own with \emph{podgorát}
  \item \emph{Recall:} say the Russian for a baking tray, then the Russian for to heat up, then say \emph{podgorát} again
\end{itemize}

\begin{tcolorbox}[breakable,colback=teal!4,colframe=teal!35!black,title={Before you move on}]
What does \ru{подгорать} mean? (\textquotedblleft{}To get burnt (of food), to catch\textquotedblright{}.) Say it once more.
\end{tcolorbox}
\section[\texorpdfstring{\ru{обжаривать} (obzhárivat)}{обжаривать (obzhárivat)}]{\ru{обжаривать} (obzhárivat) — to brown, to fry lightly}

\label{lesson:RU-C332-obzharivat}



\begin{quote}
Before the new one: say the Russian for to heat up, then the Russian for to get burnt.
\end{quote}

\subsection*{You'll want to know: \ru{обжаривать}}
\textbf{\ru{обжаривать}} (\emph{obzhárivat}) — \textquotedblleft{}to brown, to fry lightly\textquotedblright{}.

\begin{quote}
\ru{Он} \ru{обжарил} \ru{лук} \ru{на} \ru{масле}.
\end{quote}

\begin{quote}
\emph{On obzháril luk na másle.}
\end{quote}

\begin{quote}
He browned the onion in oil.
\end{quote}

Imperfective; perfective \emph{\ru{обжарить}} (\emph{obzhárit}). \emph{\ru{Жарить}} — to fry.

\begin{cousinweb}
\emph{\ru{обжаривать}} is built on \emph{\ru{об}-} \textquotedblleft{}all round\textquotedblright{} and \emph{\ru{жарить}} \textquotedblleft{}to fry\textquotedblright{}, from \emph{\ru{жар}} \textquotedblleft{}heat\textquotedblright{}, the same Slavic root as \emph{\ru{гореть}} \textquotedblleft{}to burn\textquotedblright{}.
\end{cousinweb}

\begin{grammarlens}[title={What you've built}]
One more word for telling, explaining and describing at B1.
\end{grammarlens}

\subsection*{Your turn}
\begin{itemize}
\raggedright
  \item \emph{Say it:} \emph{obzhárivat}
  \item \emph{Say it:} \emph{obzhárivat}, once more
  \item \emph{Say it:} a sentence of your own with \emph{obzhárivat}
  \item \emph{Recall:} say the Russian for to heat up, then the Russian for to get burnt, then say \emph{obzhárivat} again
\end{itemize}

\begin{tcolorbox}[breakable,colback=teal!4,colframe=teal!35!black,title={Before you move on}]
What does \ru{обжаривать} mean? (\textquotedblleft{}To brown, to fry lightly\textquotedblright{}.) Say it once more.
\end{tcolorbox}
